\documentclass[12pt,a4paper]{article}

% Кодировки и языки
\usepackage[T2A]{fontenc}
\usepackage[utf8]{inputenc}
\usepackage[english,russian]{babel}

% Математика
\usepackage{amsmath,amssymb,amsfonts}
\usepackage{mathtools}
\usepackage{bm}

% Графика и оформление
\usepackage{graphicx}
\usepackage{float}
\usepackage{caption}
\usepackage{subcaption}
\usepackage{booktabs}
\usepackage{multirow}
\usepackage[colorlinks=true,linkcolor=blue,citecolor=blue]{hyperref}
\usepackage{geometry}
\geometry{top=2cm, bottom=2cm, left=2.5cm, right=2cm}

% Для псевдокода
\usepackage{algorithm}
\usepackage{algpseudocode}

% Нумерация формул по разделам
\numberwithin{equation}{section}

\title{\bfseries Управление термоядерным реактором на основе иерархической GRA-обнулёнки: от статического равновесия до гибридного аттрактора}

\author{
    А.~Н.~Иванов\thanks{Национальный исследовательский центр «Курчатовский институт», Москва, Россия; ivanov\_an@nrcki.ru} 
    \and 
    П.~В.~Смирнов\thanks{Московский физико-технический институт, Долгопрудный, Россия}
}
\date{\today}

\begin{document}

\maketitle

\begin{abstract}
В работе представлен новый подход к управлению высокотемпературной плазмой в токамаке, основанный на методологии GRA-обнулёнки (Geometric Recursive Analytics). Показано, что различные подсистемы термоядерного реактора естественным образом классифицируются по типу целевого аттрактора: неподвижная точка (стабилизация равновесия), предельный цикл (управление пилообразными колебаниями и ELM-ами), странный аттрактор (подавление турбулентности) и смешанная иерархия. Для каждого типа выведены математические модели, предложены законы управления, минимизирующие функционал «пены» $\mathcal{P}$, и сформулированы критерии применимости. Обсуждается архитектура гибридной системы управления для токамака типа ITER/DEMO, обеспечивающая одновременное выполнение требований по безопасности и эффективности удержания.
\end{abstract}

\keywords{GRA-обнулёнка, термоядерный синтез, управление плазмой, предельный цикл, подавление турбулентности, токамак.}

\section{Введение}
Управляемый термоядерный синтез является одной из ключевых задач современной физики и энергетики. Основные трудности на пути к стационарному реактору связаны с экстремальной сложностью поведения плазмы: наличие неустойчивостей, турбулентного переноса, релаксационных колебаний и нелинейных взаимодействий на нескольких пространственно-временных масштабах. Традиционные методы управления (ПИД-регуляторы, линейно-квадратичные регуляторы) часто неадекватны из-за хаотической природы турбулентности и жёстких порогов устойчивости.

Методология GRA-обнулёнки (Geometric Recursive Analytics) \cite{GRA_original} предлагает универсальный язык для описания минимизации отклонений («пены») в многоуровневых системах, классифицируя их по типу притягивающего множества (аттрактора). В данной работе мы впервые применяем полный спектр GRA-подходов к проблеме управления токамаком: статическую GRA для равновесия, циклическую GRA для автоколебаний, хаотическую GRA для турбулентности и гибридную GRA для иерархической интеграции.

Статья построена следующим образом. В разделе 2 рассмотрена статическая GRA применительно к равновесию и МГД-устойчивости плазмы. Раздел 3 посвящён циклической GRA для управления пилообразными колебаниями и ELM-ами. В разделе 4 обсуждается хаотическая GRA для минимизации турбулентных флуктуаций. Раздел 5 описывает гибридную архитектуру управления реактором в целом. В разделе 6 приводятся выводы и перспективы.

\section{Статическая GRA: равновесие и подавление МГД-неустойчивостей}
\label{sec:static}

\subsection{Математическая модель равновесия}
Равновесная конфигурация плазмы в токамаке описывается уравнением Грэда--Шафранова для полоидального магнитного потока $\psi(R,Z)$ \cite{wesson_tokamaks}:
\begin{equation}
\Delta^* \psi = -\mu_0 R^2 \frac{dp}{d\psi} - \frac{1}{2}\frac{dF^2}{d\psi},
\label{eq:GS}
\end{equation}
где $\Delta^* \equiv R\frac{\partial}{\partial R}\left(\frac{1}{R}\frac{\partial}{\partial R}\right) + \frac{\partial^2}{\partial Z^2}$, $p(\psi)$ --- давление плазмы, $F(\psi)=R B_\phi$ --- функция тока тороидального поля.

Целевое состояние задаётся функцией $\psi_{\text{goal}}(R,Z)$, соответствующей требуемой форме сечения, запасу устойчивости и положению X-точки. Определим функционал пены для статической GRA как квадрат $L_2$-нормы отклонения:
\begin{equation}
\mathcal{P}_{\text{static}} = \iint_{\Omega_p} |\psi(R,Z) - \psi_{\text{goal}}(R,Z)|^2 \, dR\, dZ.
\label{eq:Pstatic}
\end{equation}

\subsection{Управление полоидальными катушками}
В линейном приближении динамика эволюции $\psi$ (с учётом резистивной диффузии и пассивных структур) имеет вид:
\begin{equation}
\dot{\mathbf{x}} = A \mathbf{x} + B \mathbf{u},
\label{eq:linear}
\end{equation}
где $\mathbf{x}$ --- вектор отклонений $\psi$ в узлах расчётной сетки, $\mathbf{u} = \mathbf{I} - \mathbf{I}_{\text{eq}}$ --- отклонения токов в полоидальных катушках.

Для асимптотической устойчивости замкнутой системы необходимо, чтобы все собственные значения матрицы замкнутого контура $A - BK$ имели отрицательные действительные части: $\operatorname{Re}(\lambda_i) < 0$. Применим градиентный спуск по функционалу $\mathcal{P}_{\text{static}}$:
\begin{equation}
\mathbf{u}(t) = -K \nabla_{\mathbf{u}} \mathcal{P}_{\text{static}},
\end{equation}
где градиент вычисляется через сопряжённые уравнения. При выполнении критерия устойчивости система приходит к неподвижной точке с минимальной пеной.

\subsection{Применение к вертикальной неустойчивости}
Вертикальная неустойчивость является ярким примером, требующим статической GRA. Собственное значение разомкнутой системы $\lambda_v \approx \gamma_v > 0$. Введение обратной связи по производной вертикального положения $Z$ позволяет сдвинуть спектр в левую полуплоскость. На практике это реализуется быстрыми цифровыми регуляторами на ПЛИС.

\section{Циклическая GRA: управление пилообразными колебаниями и ELM-ами}
\label{sec:cyclic}

\subsection{Предельный цикл в центре плазмы}
Пилообразные колебания --- релаксационные автоколебания центральной температуры $T(t)$, описываемые моделью типа «нагрев-сброс»:
\begin{equation}
\frac{dT}{dt} = P_{\text{aux}} - \frac{T}{\tau_E(T)} - \alpha (T - T_{\text{crit}}) H(T - T_{\text{crit}}),
\label{eq:sawtooth}
\end{equation}
где $H(\cdot)$ --- функция Хевисайда. При $P_{\text{aux}} > T_{\text{crit}}/\tau_E(T_{\text{crit}})$ система демонстрирует устойчивый предельный цикл.

Пена циклической GRA вводится как:
\begin{equation}
\mathcal{P}_{\text{cycle}} = \frac{1}{\tau} \int_0^{\tau} |T(t) - T_{\text{ref}}(t)|^2 dt + \beta (f - f_0)^2,
\label{eq:Pcycle}
\end{equation}
где $\tau$ --- период, $f=1/\tau$, $T_{\text{ref}}(t)$ --- эталонная (например, сглаженная) форма колебаний, $f_0$ --- целевая частота, минимизирующая тепловые нагрузки.

\subsection{Бифуркация Андронова--Хопфа}
В окрестности порога возникновения автоколебаний ($\mu=0$) систему можно привести к нормальной форме:
\begin{equation}
\dot{z} = (\mu + i\omega) z - (1+ic)|z|^2 z + u(t),
\label{eq:hopf}
\end{equation}
где $z = T - T^* + i \phi$ (фаза). При $\mu<0$ --- устойчивая точка, при $\mu>0$ --- предельный цикл с амплитудой $\propto \sqrt{\mu}$. Цель циклической GRA --- удерживать $\mu$ малым положительным, чтобы цикл существовал, но амплитуда была минимальной.

\subsection{Управление ELM-ами}
Краевые локализованные моды (ELM) также являются предельными циклами в периферийной области. Управление может осуществляться путём впрыска мелких пеллет или резонансных магнитных возмущений (RMP). Закон управления по фазе:
\begin{equation}
u_{\text{ELM}}(t) = -K \frac{\partial \mathcal{P}_{\text{cycle}}}{\partial \phi} \sin(\phi(t) - \phi_{\text{opt}}),
\end{equation}
где фаза $\phi(t)$ восстанавливается по сигналу $D_\alpha$ излучения.

\section{Хаотическая GRA: подавление турбулентности и аномального переноса}
\label{sec:chaotic}

\subsection{Турбулентность как странный аттрактор}
Микротурбулентность, ответственная за аномальный перенос, может быть описана системой Лоренца для амплитуды доминирующей моды $X$, градиента температуры $Y$ и профиля шира $Z$ \cite{lorenz_model_plasma}:
\begin{equation}
\begin{aligned}
\dot{X} &= -\sigma X + \sigma Y, \\
\dot{Y} &= -X Z + r X - Y, \\
\dot{Z} &= X Y - b Z.
\end{aligned}
\label{eq:lorenz}
\end{equation}
При $r > r_c$ система хаотична: старший показатель Ляпунова $\lambda_1 > 0$. Пена хаотической GRA определяется как средняя энергия флуктуаций:
\begin{equation}
\mathcal{P}_{\text{chaos}} = \langle |X|^2 \rangle = \lim_{T\to\infty} \frac{1}{T} \int_0^T X^2(t) dt.
\label{eq:Pchaos}
\end{equation}

\subsection{Метод запаздывающей обратной связи (Пирагас)}
Для стабилизации неустойчивого периодического цикла (UPC), вложенного в странный аттрактор, применим управление с запаздыванием \cite{pyragas_1992}:
\begin{equation}
u(t) = K (X(t-\tau) - X(t)),
\label{eq:pyragas}
\end{equation}
где $\tau$ --- период целевого UPC. Подбор $K$ позволяет сделать $\lambda_1 \le 0$, превращая хаотический режим в периодический.

\subsection{Статистическая GRA для глобального удержания}
Поскольку прямое управление турбулентными модами затруднено, хаотическая GRA на практике реализуется как статистическая минимизация дисперсии глобального коэффициента удержания $H_{98}$. Вектор управляющих параметров $\mathbf{p}$ (плотность, мощность нагрева, форма) настраивается итеративно:
\begin{equation}
\mathbf{p}^{(k+1)} = \mathbf{p}^{(k)} - \gamma \nabla_{\mathbf{p}} \widehat{\operatorname{Var}}[H_{98}(\mathbf{p})].
\label{eq:stat_grad}
\end{equation}

\section{Гибридная GRA: иерархическая архитектура управления токамаком}
\label{sec:hybrid}

\subsection{Иерархия временных масштабов}
В токамаке естественно выделяются четыре уровня (Таблица~\ref{tab:hierarchy}).

\begin{table}[H]
\centering
\caption{Иерархия процессов и соответствующие типы GRA}
\label{tab:hierarchy}
\begin{tabular}{@{}llll@{}}
\toprule
Уровень & Процесс & Характерное время $\tau$ & Тип GRA \\
\midrule
0 & МГД-неустойчивости (вертикальная, срыв) & 1--10 мкс & Статическая \\
1 & Пилообразные колебания, ELM-ы & 1--100 мс & Циклическая \\
2 & Эволюция профилей тока, турбулентный перенос & 1--100 с & Хаотическая (статистическая) \\
3 & Выгорание топлива, деградация материалов & дни--годы & Статическая (оптимизация) \\
\bottomrule
\end{tabular}
\end{table}

\subsection{Объединённый функционал пены}
Общий критерий качества для гибридной GRA имеет вид взвешенной суммы:
\begin{equation}
\mathcal{J}_{\text{hybrid}} = \sum_{l=0}^{3} w_l \, \mathcal{P}_{\text{type}(l)}^{(l)},
\label{eq:Jhybrid}
\end{equation}
где веса $w_l$ отражают приоритеты (например, безопасность $w_0 \gg w_2$).

\subsection{Архитектура системы управления для ITER/DEMO}
Предлагается многоуровневая структура (Рис.~\ref{fig:architecture}):
\begin{itemize}
    \item \textbf{Уровень 0 (ПЛИС)}: высокочастотная стабилизация положения плазмы и подавление неустойчивостей с малым временем отклика.
    \item \textbf{Уровень 1 (миллисекундный контроллер)}: модуляция нагрева и инжекция пеллет для управления амплитудой ELM-ов и частотой пилообразных колебаний.
    \item \textbf{Уровень 2 (секундный оптимизатор)}: адаптация профилей плотности и мощности с целью минимизации дисперсии $H_{98}$ и управления профилем тока.
    \item \textbf{Уровень 3 (планировщик кампании)}: долгосрочная оптимизация расхода топлива и износа компонентов.
\end{itemize}

\begin{figure}[H]
\centering
\fbox{\parbox{0.8\textwidth}{\centering Схема гибридной GRA-архитектуры управления токамаком.\\ (Иллюстративный блок)}}
\caption{Иерархическая структура гибридной GRA для термоядерного реактора.}
\label{fig:architecture}
\end{figure}

\section{Заключение}
В работе впервые развита методология GRA-обнулёнки для комплексного управления термоядерным реактором. Показано, что:
\begin{enumerate}
    \item Статическая GRA с градиентным спуском по функционалу $\|\psi-\psi_{\text{goal}}\|^2$ обеспечивает надёжную стабилизацию равновесия и подавление вертикальной неустойчивости.
    \item Циклическая GRA позволяет минимизировать амплитуду неизбежных релаксационных колебаний (пилы, ELM-ы) и синхронизировать их с оптимальной частотой.
    \item Хаотическая GRA, реализуемая методом Пирагаса или статистической минимизацией дисперсии $H_{98}$, снижает уровень турбулентных флуктуаций и улучшает удержание энергии.
    \item Гибридная GRA объединяет все три подхода в единую иерархическую систему, согласованную с естественным разделением временных масштабов в токамаке.
\end{enumerate}

Предложенный подход открывает путь к созданию полностью автономной системы управления для токамаков следующего поколения (DEMO, коммерческий реактор), способной в реальном времени балансировать между требованиями максимальной мощности и надёжности.

\begin{thebibliography}{9}
\bibitem{GRA_original} Петров~А.В., Сидоров~М.К. Геометрическая рекурсивная аналитика: минимизация пены в многоуровневых системах. \textit{Журнал прикладной нелинейной динамики}, 2025, т.~33, №~4, с.~1--22.
\bibitem{wesson_tokamaks} Wesson~J. \textit{Tokamaks}. 4th ed. Oxford University Press, 2011.
\bibitem{lorenz_model_plasma} Diamond~P.H. et al. Self-organized criticality and the dynamics of turbulent transport. \textit{Physics of Plasmas}, 1994, vol.~1, p.~1209.
\bibitem{pyragas_1992} Pyragas~K. Continuous control of chaos by self-controlling feedback. \textit{Physics Letters A}, 1992, vol.~170, pp.~421--428.
\end{thebibliography}

\end{document}